% GENERATED FILE. Edit canonical lessons, then run npm run generate:books.
% canonical-source-hash: fnv1a64:506310b71271abe1
% canonical-lessons: SA-C440-himapatah, SA-C440-kantha, SA-C440-viliyate, SA-C440-sankucati, SA-C440-pravarah, SA-R440-review-434-440

\chapter{Snow and Cold}
\label{ch:sa-snow-and-cold}
\hlchaptermodality{440}

\begin{chapteropening}
\textbf{By the end of this chapter:} \emph{I can talk about snowfall, melting snow and keeping warm.}

Complete the last lesson of chapter 440: I can talk about snowfall, melting snow and keeping warm.
\end{chapteropening}

\section[\texorpdfstring{\sk{हिमपातः} (himapātaḥ)}{हिमपातः (himapātaḥ)}]{\sk{हिमपातः} (himapātaḥ) — a snowfall}

\label{lesson:SA-C440-himapatah}



\begin{quote}
Before the new one: say the Sanskrit for sets swaying, then the Sanskrit for a gentle breeze.
\end{quote}

\subsection*{You'll want to know: \sk{हिमपातः}}
\textbf{\sk{हिमपातः}} — \emph{himapātaḥ} — \textquotedblleft{}a snowfall\textquotedblright{}.

\begin{quote}
\sk{पर्वते} \sk{रात्रौ} \sk{हिमपातः} \sk{भवति।}
\end{quote}

\begin{quote}
\emph{parvate rātrau himapātaḥ bhavati.}
\end{quote}

\begin{quote}
There is a snowfall on the mountain at night.
\end{quote}

A masculine noun.

\begin{cousinweb}
\textbf{\sk{हिमपातः}} (\emph{himapātaḥ}) is built from \emph{hima}, \textquotedblleft{}snow\textquotedblright{}, and \emph{pāta}, \textquotedblleft{}a fall\textquotedblright{}, from the root √pat, \textquotedblleft{}to fall\textquotedblright{}, which also gives \emph{patati}, \textquotedblleft{}falls\textquotedblright{}; MW gives the compound as a fall of snow (MW).
\end{cousinweb}

\begin{grammarlens}[title={What you've built}]
One more word for getting around.
\end{grammarlens}

\subsection*{Your turn}
\begin{itemize}
\raggedright
  \item \emph{Say it:} \emph{himapātaḥ}
  \item \emph{Say it:} \emph{himapātaḥ}, once more
  \item \emph{Say it:} \emph{parvate rātrau himapātaḥ bhavati.} — There is a snowfall on the mountain at night.
  \item \emph{Recall:} say the Sanskrit for sets swaying, then the Sanskrit for a gentle breeze, then say \emph{himapātaḥ} again
\end{itemize}

\begin{tcolorbox}[breakable,colback=teal!4,colframe=teal!35!black,title={Before you move on}]
What does \sk{हिमपातः} mean? (\textquotedblleft{}A snowfall\textquotedblright{}.) Say it once more.
\end{tcolorbox}
\section[\texorpdfstring{\sk{कन्था} (kanthā)}{कन्था (kanthā)}]{\sk{कन्था} (kanthā) — a patched garment; a quilt of rags}

\label{lesson:SA-C440-kantha}



\begin{quote}
Before the new one: say the Sanskrit for a gentle breeze, then the Sanskrit for a snowfall.
\end{quote}

\subsection*{You'll want to know: \sk{कन्था}}
\textbf{\sk{कन्था}} — \emph{kanthā} — \textquotedblleft{}a patched garment; a quilt of rags\textquotedblright{}.

\begin{quote}
\sk{शीते} \sk{भिक्षुः} \sk{कन्थां} \sk{धारयति।}
\end{quote}

\begin{quote}
\emph{śīte bhikṣuḥ kanthāṃ dhārayati.}
\end{quote}

\begin{quote}
In the cold the monk wears a patched garment.
\end{quote}

A feminine noun.

\begin{cousinweb}
\textbf{\sk{कन्था}} (\emph{kanthā}) is the word MW gives for a rag or a patched garment, especially one that ascetics wear; MW gives no derivation, so its origin is uncertain; it lives on as Hindi \emph{kathrī}, \textquotedblleft{}a quilt of patches\textquotedblright{} (MW; CDIAL 2721).
\end{cousinweb}

\begin{grammarlens}[title={What you've built}]
One more word for getting around.
\end{grammarlens}

\subsection*{Your turn}
\begin{itemize}
\raggedright
  \item \emph{Say it:} \emph{kanthā}
  \item \emph{Say it:} \emph{kanthā}, once more
  \item \emph{Say it:} \emph{śīte bhikṣuḥ kanthāṃ dhārayati.} — In the cold the monk wears a patched garment.
  \item \emph{Recall:} say the Sanskrit for a gentle breeze, then the Sanskrit for a snowfall, then say \emph{kanthā} again
\end{itemize}

\begin{tcolorbox}[breakable,colback=teal!4,colframe=teal!35!black,title={Before you move on}]
What does \sk{कन्था} mean? (\textquotedblleft{}A patched garment; a quilt of rags\textquotedblright{}.) Say it once more.
\end{tcolorbox}
\section[\texorpdfstring{\sk{विलीयते} (vilīyate)}{विलीयते (vilīyate)}]{\sk{विलीयते} (vilīyate) — melts}

\label{lesson:SA-C440-viliyate}



\begin{quote}
Before the new one: say the Sanskrit for a snowfall, then the Sanskrit for a patched garment.
\end{quote}

\subsection*{You'll want to know: \sk{विलीयते}}
\textbf{\sk{विलीयते}} — \emph{vilīyate} — \textquotedblleft{}melts\textquotedblright{}.

\begin{quote}
\sk{ग्रीष्मे} \sk{पर्वते} \sk{हिमं} \sk{विलीयते।}
\end{quote}

\begin{quote}
\emph{grīṣme parvate himaṃ vilīyate.}
\end{quote}

\begin{quote}
In summer the snow on the mountain melts.
\end{quote}

Past (it …): \sk{हिमं} \sk{विलीनवत्}, \emph{himaṃ vilīnavat}.

\begin{cousinweb}
\textbf{\sk{विलीयते}} (\emph{vilīyate}) is built from \emph{vi}, \textquotedblleft{}apart\textquotedblright{}, and the root √lī, \textquotedblleft{}to cling, to dissolve\textquotedblright{}; MW gives \emph{vi-lī} as to be dissolved, to melt, with the participle \emph{vilīna}, \textquotedblleft{}melted\textquotedblright{} (MW).
\end{cousinweb}

\begin{grammarlens}[title={What you've built}]
One more word for getting around.
\end{grammarlens}

\subsection*{Your turn}
\begin{itemize}
\raggedright
  \item \emph{Say it:} \emph{vilīyate}
  \item \emph{Say it:} \emph{vilīyate}, once more
  \item \emph{Say it:} \emph{grīṣme parvate himaṃ vilīyate.} — In summer the snow on the mountain melts.
  \item \emph{Recall:} say the Sanskrit for a snowfall, then the Sanskrit for a patched garment, then say \emph{vilīyate} again
\end{itemize}

\begin{tcolorbox}[breakable,colback=teal!4,colframe=teal!35!black,title={Before you move on}]
What does \sk{विलीयते} mean? (\textquotedblleft{}Melts\textquotedblright{}.) Say it once more.
\end{tcolorbox}
\section[\texorpdfstring{\sk{संकुचति} (saṅkucati)}{संकुचति (saṅkucati)}]{\sk{संकुचति} (saṅkucati) — huddles up; shrinks}

\label{lesson:SA-C440-sankucati}



\begin{quote}
Before the new one: say the Sanskrit for a patched garment, then the Sanskrit for melts.
\end{quote}

\subsection*{You'll want to know: \sk{संकुचति}}
\textbf{\sk{संकुचति}} — \emph{saṅkucati} — \textquotedblleft{}huddles up; shrinks\textquotedblright{}.

\begin{quote}
\sk{शीते} \sk{बालः} \sk{कम्बले} \sk{संकुचति।}
\end{quote}

\begin{quote}
\emph{śīte bālaḥ kambale saṅkucati.}
\end{quote}

\begin{quote}
In the cold the boy huddles up in a blanket.
\end{quote}

Past (he …): \sk{सः} \sk{कम्बले} \sk{संकुचितवान्}, \emph{saḥ kambale saṅkucitavān}.

\begin{cousinweb}
\textbf{\sk{संकुचति}} (\emph{saṅkucati}) is built from \emph{sam}, \textquotedblleft{}together\textquotedblright{}, and the root √kuc, \textquotedblleft{}to contract, to bend\textquotedblright{}; MW gives \emph{saṃ-kuc} as to contract or shrink, and its participle \emph{saṃkucita} as crouching, cowering (MW).
\end{cousinweb}

\begin{grammarlens}[title={What you've built}]
One more word for getting around.
\end{grammarlens}

\subsection*{Your turn}
\begin{itemize}
\raggedright
  \item \emph{Say it:} \emph{saṅkucati}
  \item \emph{Say it:} \emph{saṅkucati}, once more
  \item \emph{Say it:} \emph{śīte bālaḥ kambale saṅkucati.} — In the cold the boy huddles up in a blanket.
  \item \emph{Recall:} say the Sanskrit for a patched garment, then the Sanskrit for melts, then say \emph{saṅkucati} again
\end{itemize}

\begin{tcolorbox}[breakable,colback=teal!4,colframe=teal!35!black,title={Before you move on}]
What does \sk{संकुचति} mean? (\textquotedblleft{}Huddles up; shrinks\textquotedblright{}.) Say it once more.
\end{tcolorbox}
\section[\texorpdfstring{\sk{प्रावारः} (prāvāraḥ)}{प्रावारः (prāvāraḥ)}]{\sk{प्रावारः} (prāvāraḥ) — a cloak}

\label{lesson:SA-C440-pravarah}



\begin{quote}
Before the new one: say the Sanskrit for melts, then the Sanskrit for huddles up.
\end{quote}

\subsection*{You'll want to know: \sk{प्रावारः}}
\textbf{\sk{प्रावारः}} — \emph{prāvāraḥ} — \textquotedblleft{}a cloak\textquotedblright{}.

\begin{quote}
\sk{शीते} \sk{वृद्धः} \sk{प्रावारं} \sk{धारयति।}
\end{quote}

\begin{quote}
\emph{śīte vṛddhaḥ prāvāraṃ dhārayati.}
\end{quote}

\begin{quote}
In the cold the old man wears a cloak.
\end{quote}

A masculine noun.

\begin{cousinweb}
\textbf{\sk{प्रावारः}} (\emph{prāvāraḥ}) is built from \emph{prā}, a form of \emph{pra}, and the root √vṛ, \textquotedblleft{}to cover\textquotedblright{}; MW gives \emph{prāvāra} as an outer garment, a mantle, a cloak (MW).
\end{cousinweb}

\begin{grammarlens}[title={What you've built}]
One more word for getting around.
\end{grammarlens}

\subsection*{Your turn}
\begin{itemize}
\raggedright
  \item \emph{Say it:} \emph{prāvāraḥ}
  \item \emph{Say it:} \emph{prāvāraḥ}, once more
  \item \emph{Say it:} \emph{śīte vṛddhaḥ prāvāraṃ dhārayati.} — In the cold the old man wears a cloak.
  \item \emph{Recall:} say the Sanskrit for melts, then the Sanskrit for huddles up, then say \emph{prāvāraḥ} again
\end{itemize}

\begin{tcolorbox}[breakable,colback=teal!4,colframe=teal!35!black,title={Before you move on}]
What does \sk{प्रावारः} mean? (\textquotedblleft{}A cloak\textquotedblright{}.) Say it once more.
\end{tcolorbox}
\section[(dialogue)]{Looking back: from Loans and Deposits to Snow and Cold}

\label{lesson:SA-R440-review-434-440}



\begin{quote}
Say the Sanskrit for a creditor and a debtor.
\end{quote}

\subsection*{Your turn}
Say these aloud:

\begin{itemize}
\raggedright
  \item a creditor, a debtor, money lent at interest, a deposit, pays back
  \item a cowrie shell, a gold piece, a dinar, barter, a storehouse
  \item a bank of cloud, rainfall, starts to pour, gets soaked, slippery
  \item the rainy season, a downpour, drizzle, a drop, drips
  \item sunshine, the hot season, scorching, sweat, blazes
  \item a breeze, a tempest, violent, sets swaying, a gentle breeze
  \item a snowfall, a patched garment, melts, huddles up, a cloak
\end{itemize}

\begin{tcolorbox}[breakable,colback=teal!4,colframe=teal!35!black,title={Before you move on}]
A creditor? (\textbf{\sk{उत्तमर्णः}}, \emph{uttamarṇaḥ}.) A debtor? (\textbf{\sk{अधमर्णः}}, \emph{adhamarṇaḥ}.) Money lent at interest? (\textbf{\sk{कुसीदम्}}, \emph{kusīdam}.) A deposit? (\textbf{\sk{निक्षेपः}}, \emph{nikṣepaḥ}.) Pays back? (\textbf{\sk{निर्यातयति}}, \emph{niryātayati}.) A cowrie shell? (\textbf{\sk{कपर्दकः}}, \emph{kapardakaḥ}.) A gold piece? (\textbf{\sk{निष्कः}}, \emph{niṣkaḥ}.) A dinar? (\textbf{\sk{दीनारः}}, \emph{dīnāraḥ}.) Barter? (\textbf{\sk{विनिमयः}}, \emph{vinimayaḥ}.) A storehouse? (\textbf{\sk{भाण्डागारम्}}, \emph{bhāṇḍāgāram}.) A bank of cloud? (\textbf{\sk{अभ्रम्}}, \emph{abhram}.) Rainfall? (\textbf{\sk{वृष्टिः}}, \emph{vṛṣṭiḥ}.) Starts to pour? (\textbf{\sk{प्रवर्षति}}, \emph{pravarṣati}.) Gets soaked? (\textbf{\sk{क्लिद्यति}}, \emph{klidyati}.) Slippery? (\textbf{\sk{पिच्छिलः}}, \emph{picchilaḥ}.) The rainy season? (\textbf{\sk{प्रावृट्}}, \emph{prāvṛṭ}.) A downpour? (\textbf{\sk{आसारः}}, \emph{āsāraḥ}.) Drizzle? (\textbf{\sk{शीकरः}}, \emph{śīkaraḥ}.) A drop? (\textbf{\sk{बिन्दुः}}, \emph{binduḥ}.) Drips? (\textbf{\sk{गलति}}, \emph{galati}.) Sunshine? (\textbf{\sk{आतपः}}, \emph{ātapaḥ}.) The hot season? (\textbf{\sk{निदाघः}}, \emph{nidāghaḥ}.) Scorching? (\textbf{\sk{प्रखरः}}, \emph{prakharaḥ}.) Sweat? (\textbf{\sk{स्वेदः}}, \emph{svedaḥ}.) Blazes? (\textbf{\sk{दीप्यते}}, \emph{dīpyate}.) A breeze? (\textbf{\sk{पवनः}}, \emph{pavanaḥ}.) A tempest? (\textbf{\sk{प्रभञ्जनः}}, \emph{prabhañjanaḥ}.) Violent? (\textbf{\sk{प्रबलः}}, \emph{prabalaḥ}.) Sets swaying? (\textbf{\sk{आन्दोलयति}}, \emph{āndolayati}.) A gentle breeze? (\textbf{\sk{मन्दानिलः}}, \emph{mandānilaḥ}.) A snowfall? (\textbf{\sk{हिमपातः}}, \emph{himapātaḥ}.) A patched garment? (\textbf{\sk{कन्था}}, \emph{kanthā}.) Melts? (\textbf{\sk{विलीयते}}, \emph{vilīyate}.) Huddles up? (\textbf{\sk{संकुचति}}, \emph{saṅkucati}.) A cloak? (\textbf{\sk{प्रावारः}}, \emph{prāvāraḥ}.)
\end{tcolorbox}
